\documentclass[a4paper,12pt]{article}

% --- Use XeLaTeX font support ---

\usepackage{fontspec}
\setmainfont{TeX Gyre Termes} % nice Times-like font, installed by texlive-core
\setsansfont{TeX Gyre Heros} % optional sans-serif
\setmonofont{Fira Code}       % optional monospaced font
% --- Math packages ---
\usepackage{amsmath, amssymb, amsthm}
\usepackage{mathtools}

% --- Page layout ---
\usepackage[margin=1in]{geometry}

% --- Better lists ---
\usepackage{enumitem}
% --- Hyperlinks for references ---
\usepackage[colorlinks=true, linkcolor=blue, urlcolor=blue]{hyperref}

% Header: fancy layout like the rovided image
\usepackage{fancyhdr}
\setlength{\headheight}{26pt} % increase if fancyhdr warns
\pagestyle{fancy}
\fancyhf{} % clear header/footer
\rhead{\begin{tabular}{r} Mt.: 3052737\end{tabular}}
\chead{\begin{tabular}{c}Lösungen EA 2\end{tabular}}
\lhead{\begin{tabular}{l}61521 Einführung in die \\ Numerische Mathematik SS26 \end{tabular}}
\renewcommand{\headrulewidth}{0.8pt}
\fancyfoot[C]{- \thepage\ -}
% Ensure the plain page style (used by \maketitle) uses the same header

% --- Line spacing ---
\usepackage{setspace} % provides \doublespacing and \onehalfspacing
\usepackage{booktabs}
% --- Optional solutions environment ---
\newenvironment{solution}{
  \par\noindent\textbf{Solution:}\quad
}{\hfill$\blacksquare$\par\vspace{1em}}

% % --- Title info ---
% \title{Aufgabe \#1}
% \author{Your Name}
% \date{\today}
% 
\begin{document}
\raggedright
% Enable double (2.0) line spacing for the document body
\doublespacing
\noindent\textbf{Aufgabe 4} \\
%\noindent\textbf{(a)} \\
Gegeben:
\begin{itemize}
    \item $f : \mathbb{R}^n \to \mathbb{R}^n$
    \item $f(x^*) = 0$
    \item $f'(x^*)$ invertierbar
    \item $f'$ lokal Lipschitz-stetig
\end{itemize}

Iteration:
\[
x^{k+1} = x^k + d^k
\]

mit
\[
\|f'(x^k)d^k + f(x^k)\|_\infty \le \|f(x^k)\|_\infty^\tau, \quad \tau \in (1, 2].
\]
Setze 
$$e_k:=x^k - x^*.$$
Schreibe
$$f'(x^k)d^k=-f(x^k)+r_k$$
mit
$$||r_k|| \leq ||f(x^k)||_{\infty}^{\tau}$$
Da $f'$ Lipschitz-stetig ist
$$
f(x^k)=f'(x^*)(x^k-x^*)+R_k=f'(x^*)e_k +R_k
$$
mit
$$||R_k|| \leq C||e_k||^2.
$$
Da $f'$ stetig ist:
$$
f'(x^k) = f'(x^*) + O(\| e_k \|).
$$
Einsetzen:
$$
f'(x^k)d^k = -f'(x^*) e_k - R_k + r_k.
$$
Gesucht wird ein Ausdruck für $e_{k+1}$: 
$$
\boxed{e_{k+1} = x^{k+1} - x^* = x^k + d^k - x^* = e_k + d^k.}
$$
Auflösen nach $d_k$ dazu
multipliziere mit $(f'(x^k))^{-1}$:
$$
d^k = -(f'(x^k))^{-1} f'(x^*) e_k - (f'(x^k))^{-1} R_k + (f'(x^k))^{-1} r_k.
$$
Einsetzen in $e_{k+1}$
$$
e_{k+1} = e_k - (f'(x^k))^{-1} f'(x^*) e_k - (f'(x^k))^{-1} R_k + (f'(x^k))^{-1} r_k.
$$
Der Ausdruck
$$
I - (f'(x^k))^{-1} f'(x^*)
$$
ist klein:
$= O(\| e_k \|) \quad \text{(wegen Stetigkeit von $f'$)}$.\\
Abschätzen der Terme:
\begin{itemize}
  \item Linearer Term
$$
\left\| e_k - (f'(x^k))^{-1} f'(x^*) e_k \right\| \le C \| e_k \|^2.
$$
\item  Restterm
$$
\| (f'(x^k))^{-1} R_k \| \le C \| e_k \|^2.
$$
\item Inexaktheit
$$
\| (f'(x^k))^{-1} r_k \| \le C \| f(x^k) \|^\tau.
$$
Und $\| f(x^k) \| \le C \| e_k \|$, also
$$
\| (f'(x^k))^{-1} r_k \| \le C \| e_k \|^\tau.
$$
\end{itemize}
Gesamtschätzung:
$$
\| e_{k+1} \| \le C \| e_k \|^2 + C \| e_k \|^\tau.
$$
Da $\tau \in (1,2]$: \\
Für kleine $\| e_k \|$ gilt $\| e_k \|^\tau \gg \| e_k \|^2$ (falls $\tau < 2$). Somit
$$
\| e_{k+1} \| \le C \| e_k \|^\tau.
$$
as ist genau die Situation aus Teil Aufgabe1(a):
$$
\| x_{k+1} - x^* \| \le M \| x_k - x^* \|^\tau
$$
$\Rightarrow$ Q-Ordnung mindestens $\tau$.

Wie in Aufgabe(a)(i): Die rechte Seite ist klein $\Rightarrow$ Kontraktion $\Rightarrow$ Fixpunktsatz (lokal) $\Rightarrow$ Konvergenz.

\begin{itemize}
  \item Verfahren ist lokal wohldefiniert.
  \item Folge konvergiert gegen $x^*$.
  \item Konvergenzordnung:
    \[
    \boxed{\text{Q-Ordnung } \ge \tau}
    \]
\end{itemize}


\end{document}
